\begin{center}
    \Large\textbf{DAFTAR SIMBOL}
\end{center}

\renewcommand{\arraystretch}{1.2}

\setlength{\LTleft}{0pt}
\setlength{\LTright}{0pt}

\begin{longtable}{p{2.8cm} p{10.8cm}}

\textbf{Simbol} & \textbf{Keterangan} \\ \hline

\multicolumn{2}{l}{\bfseries Simbol Data dan Rolling Window} \\ \hline
$y_t$ & Nilai observasi deret waktu pada waktu ke-$t$ \\
$W_t$ & Himpunan observasi deret waktu pada rolling window ke-$t$ \\
$L$ & Panjang rolling window \\
$T$ & Jumlah total observasi deret waktu \\

\multicolumn{2}{l}{\bfseries Simbol Kompleksitas dan Ruang Parameter} \\ \hline
$c_t$ & Indeks kompleksitas data pada rolling window ke-$t$ yang dihasilkan melalui mekanisme Model Complexity Assessment (MCA) \\
$\tau_1,\tau_2$ & Threshold klasifikasi kompleksitas data pada mekanisme MCA \\
$\Omega$ & Ruang parameter global model \\
$\Theta_t$ & Ruang parameter adaptif yang digunakan dalam proses optimasi pada rolling window ke-$t$ \\
$\Theta_{t+1}$ & Ruang parameter adaptif pada window berikutnya yang diperbarui berdasarkan informasi diagnostik residual $\mathcal{D}(e_t)$ \\
$g(\cdot)$ & Mekanisme pembaruan ruang parameter adaptif berbasis diagnostik residual \\
MCA & Model Complexity Assessment \\

\multicolumn{2}{l}{\bfseries Simbol Optimasi dan Parameter Model} \\ \hline
$\theta$ & Vektor parameter model \\
$\theta_t^*$ & Parameter optimal model pada rolling window ke-$t$ \\
$\mathcal{F}$ & Operator recursive adaptive optimization untuk evolusi parameter model \\
$L_t(\theta)$ & Fungsi loss yang mengukur kesalahan prediksi model pada rolling window ke-$t$ \\
$f_t(\theta)$ & Fungsi objektif optimasi dengan komponen penalti constraint model \\
$\lambda$ & Koefisien penalti pada fungsi objektif \\
$\phi_i$ & Koefisien autoregressive orde ke-$i$ \\
$\psi_i$ & Koefisien moving average orde ke-$i$ \\
$\sigma^2$ & Variansi residual model \\
$p,d,q$ & Orde model ARIMA \\

\multicolumn{2}{l}{\bfseries Simbol Residual dan Diagnostik} \\ \hline
$e_t$ & Residual model pada waktu ke-$t$ \\
$\mathcal{D}(e_t)$ & Informasi diagnostik residual yang mencakup autokorelasi, heteroskedastisitas, skewness, dan kurtosis residual \\
$\hat{\rho}_k$ & Koefisien autokorelasi residual pada lag ke-$k$ \\
$Q$ & Statistik uji Ljung--Box \\

\multicolumn{2}{l}{\bfseries Simbol Evaluasi Model} \\ \hline
RMSE & Root Mean Squared Error sebagai ukuran rata-rata kuadrat error prediksi \\
MAE & Mean Absolute Error sebagai ukuran rata-rata absolut error prediksi \\
MAPE & Mean Absolute Percentage Error sebagai ukuran persentase error prediksi \\
$R^2$ & Koefisien determinasi \\

\end{longtable}
